\documentclass[UTF8,11pt]{ctexart}
\usepackage[a4paper,top=20mm,bottom=20mm,left=23mm,right=23mm]{geometry}
\usepackage{amsmath,amssymb,mathtools}
\usepackage{booktabs,tabularx,array}
\usepackage{enumitem}
\usepackage{fancyhdr}
\usepackage{microtype}

\pagestyle{fancy}
\fancyhf{}
\fancyhead[L]{2026 年全国中学生数学奥林匹克竞赛（预赛）一试（A卷）}
\fancyhead[R]{参考解答}
\fancyfoot[C]{\thepage}
\renewcommand{\headrulewidth}{0.4pt}
\setlength{\headheight}{14pt}
\setlength{\parindent}{2em}
\setlength{\parskip}{2pt}
\linespread{1.28}
\setlist[itemize]{leftmargin=2.2em,itemsep=2pt,topsep=3pt}
\newcommand{\ans}[1]{\boxed{#1}}

\begin{document}

\begin{center}
  {\LARGE\bfseries 2026 年全国中学生数学奥林匹克竞赛（预赛）\par}
  \vspace{3pt}
  {\Large\bfseries 一试（A卷）参考解答}
\end{center}

\section*{填空题答案}

\begin{center}
\renewcommand{\arraystretch}{1.45}
\begin{tabular}{c>{$}c<{$}@{\qquad}c>{$}c<{$}}
\toprule
题号 & \multicolumn{1}{c}{答案} & 题号 & \multicolumn{1}{c}{答案}\\
\midrule
1 & 4 & 5 & \dfrac{3\sqrt{10}}5\\
2 & 26 & 6 & -\sqrt{\dfrac{2}{\sqrt3}}+\mathrm i\\
3 & \dfrac{22}{117} & 7 & \dfrac{631}{435}\\
4 & 19 & 8 & \dfrac{227}{231}\\
\bottomrule
\end{tabular}
\end{center}

\section*{1. 平面向量}

由
\[
  \vec x\cdot\vec y=2=|\vec x|\,|\vec y|
\]
及柯西不等式的取等条件，$\vec x,\vec y$ 同向。又 $|\vec y|=2|\vec x|$，故
$\vec y=2\vec x$。因此
\[
  \vec y\cdot\vec z=2\vec x\cdot\vec z=\ans{4}.
\]

\section*{2. 等差数列}

设 $a_1=a$，公差为 $d\ne0$，则
\[
  a_2=a+d,\qquad a_4=a+3d,\qquad a_8=a+7d.
\]
因为 $a_1,2a_2,4a_4,ta_8$ 依次成等差数列，所以前三项满足
\[
  2\cdot 2(a+d)=a+4(a+3d).
\]
整理得 $a=-8d$。于是前三项为
\[
  -8d,\quad -14d,\quad -20d,
\]
其公差为 $-6d$，故第四项应为 $-26d$。另一方面，$a_8=a+7d=-d$，从而
\[
  t(-d)=-26d.
\]
由 $d\ne0$ 得
\[
  \ans{t=26}.
\]

\section*{3. 三角函数}

令 $u=\tan A$，$v=\tan B$。题设两式分别给出
\[
  u-\frac1v=9,\qquad v-\frac1u=13,
\]
即
\[
  uv-1=9v=13u.
\]
由于 $A+B=\pi-C$，
\[
  \tan C=-\tan(A+B)=\frac{u+v}{uv-1}
  =\frac{u}{uv-1}+\frac{v}{uv-1}
  =\frac1{13}+\frac19
  =\ans{\frac{22}{117}}.
\]

\section*{4. 集合中的完全平方数}

由定义
\[
  S_n=\{20,21,\ldots,20+n\}
  \cup\{26,28,\ldots,26+2n\}.
\]
当 $1\leqslant n\leqslant4$ 时，$S_n$ 中没有完全平方数。当 $n=5$ 时，
$25=20+5$ 与 $36=26+2\times5$ 同时出现；在 $5\leqslant n\leqslant18$ 时，
$S_n$ 中的完全平方数始终恰为 $25,36$。当 $n=19$ 时，新增
\[
  64=26+2\times19,
\]
此时完全平方数共有三个，个数第一次成为奇数。因此
\[
  \ans{n=19}.
\]

\section*{5. 双曲线与平行线}

将 $P(1,0)$ 代入双曲线方程，得 $a^2=1$；再代入 $Q(7,12)$，得
\[
  49-\frac{144}{b^2}=1,
\]
故 $b^2=3$。于是 $c^2=a^2+b^2=4$，右焦点为 $F(2,0)$。

设 $M$ 为 $PQ$ 的中点，则 $M=(4,6)$。两条不重合的平行线分别经过 $P,Q$，且
$F$ 到两线的距离相等，所以 $F$ 位于这两条平行线的中间平行线上。该中间平行线经过
$M$，因而两条直线均平行于 $FM$。直线 $FM$ 的斜率为 $3$，故
\[
  l_1:3x-y-3=0,\qquad l_2:3x-y-9=0.
\]
因此两平行线间的距离为
\[
  \frac{|9-3|}{\sqrt{3^2+(-1)^2}}
  =\ans{\frac{3\sqrt{10}}5}.
\]

\section*{6. 复数方程}

方程为
\[
  \bigl||z|+z\bigr|+z=\mathrm i.
\]
第一项是非负实数。比较实部与虚部，可设
\[
  z=-t+\mathrm i,\qquad t\geqslant0,
\]
且 $\bigl||z|+z\bigr|=t$。又 $|z|=\sqrt{t^2+1}$，所以
\[
  \left|\sqrt{t^2+1}-t+\mathrm i\right|=t.
\]
两边平方并整理，得
\[
  t^2+2=2t\sqrt{t^2+1}.
\]
再次平方可得
\[
  (t^2+2)^2=4t^2(t^2+1)
  \quad\Longrightarrow\quad 3t^4=4.
\]
因 $t\geqslant0$，有
\[
  t=\sqrt{\frac2{\sqrt3}}.
\]
该值满足平方前的方程，故
\[
  \ans{z=-\sqrt{\frac2{\sqrt3}}+\mathrm i}.
\]

\section*{7. 小数部分的最大值}

当 $m=n$ 时，题中式子为 $0$。由对称性，不妨设 $m>n$，并令
\[
  t=\frac mn>1,\qquad k=\lfloor t\rfloor.
\]
因为 $0<n/m<1$，所以
\[
  \left\{\frac mn\right\}+\left\{\frac nm\right\}
  =t-k+\frac1t.
\]
函数 $t+1/t$ 在 $t>1$ 上严格递增。

若 $k\geqslant2$，则 $k\leqslant t<k+1$，从而
\[
  t-k+\frac1t<1+\frac1{k+1}\leqslant\frac43.
\]

若 $k=1$，则 $n<m<2n$。当 $n\leqslant15$ 时，
\[
  \frac mn\leqslant\frac{2n-1}{n}\leqslant\frac{29}{15};
\]
当 $n\geqslant16$ 时，由 $m\leqslant30$，
\[
  \frac mn\leqslant\frac{30}{16}<\frac{29}{15}.
\]
所以此时 $t$ 的最大值在 $(m,n)=(29,15)$ 处取得。相应的值为
\[
  \frac{14}{15}+\frac{15}{29}=\frac{631}{435}>\frac43.
\]
故所求最大值为
\[
  \ans{\frac{631}{435}}.
\]

\section*{8. 正方体染色概率}

从 $12$ 条棱中任选 $6$ 条染红，共有
\[
  \binom{12}{6}=924
\]
种等可能方案。以下计算反面事件：六个面均不恰含两条红棱。

先说明，在反面事件中，一个面的红棱数不可能是 $4$。若某面四棱全红，则其余仅有
两条红棱。与该面相邻的四个面各已经含有一条红棱；要使它们都不恰含两条红棱，余下
两条红棱在每个相邻面中必须同时出现或同时不出现。但该面之外的任意两条不同棱，在
这四个相邻面中的归属集合都不相同，故必有一个相邻面只含余下两棱中的一条，矛盾。

将红棱与未染红的棱互换，反面事件保持不变，而一个没有红棱的面会变成四棱全红。
所以一个面的红棱数也不可能是 $0$。于是反面事件中，每个面的红棱数只能是 $1$ 或
$3$。每条红棱属于两个面，六条红棱在各面计数之和为 $12$。设有 $r$ 个面含三条红棱，
则
\[
  3r+(6-r)=12,
\]
从而 $r=3$。

这三个含三条红棱的面中不能有两个是相对面。否则全部六条红棱都在这两个相对面的
八条棱中，另外四个面各至多含两条红棱，而它们的红棱计数之和为 $6$；若均不含两条，
每面至多含一条，总和至多为 $4$，矛盾。因此三个含三条红棱的面两两相邻，交于同一
顶点 $V$。顶点 $V$ 有 $8$ 种选择。

固定 $V$，记其对顶点为 $W$。将十二条棱分成三类：与 $V$ 相连的三条棱中有 $X$ 条
红棱，与 $W$ 相连的三条棱中有 $Z$ 条红棱，其余六条棱中有 $H$ 条红棱。分别对经过
$V$、经过 $W$ 的三个面计算红棱数之和，得到
\[
  2X+H=9,\qquad 2Z+H=3.
\]
故 $X-Z=3$。由 $0\leqslant X,Z\leqslant3$，只能有
\[
  X=3,\qquad Z=0,\qquad H=3.
\]
其余六条棱首尾相接构成一个六边形。每一对相邻棱中必须恰有一条红棱，故六边形上的
棱只能红、白交替，共有两种染法。反面事件因此共有
\[
  8\times2=16
\]
种。所求概率为
\[
  1-\frac{16}{924}=\ans{\frac{227}{231}}.
\]

\section*{9. 实数 $C$ 的最小值}

原方程组为
\[
  \begin{cases}
    2^x+\log_2y=C,\\
    2^{-x}-\log_4y=2C.
  \end{cases}
\]
将第二式乘以 $2$ 后与第一式相加，消去对数项，得到
\[
  2^x+2\cdot2^{-x}=5C.
\]
由基本不等式
\[
  5C\geqslant2\sqrt{2^x\cdot2\cdot2^{-x}}=2\sqrt2,
\]
所以 $C\geqslant2\sqrt2/5$。等号成立当且仅当
\[
  2^x=2\cdot2^{-x},
\]
即 $x=1/2$。此时由第一式取
\[
  \log_2y=\frac{2\sqrt2}{5}-\sqrt2=-\frac{3\sqrt2}{5},
  \qquad y=2^{-3\sqrt2/5}>0,
\]
原方程组确实成立。因此
\[
  \ans{C_{\min}=\frac{2\sqrt2}{5}}.
\]

\section*{10. 小钢球个数的最大值}

以圆柱轴线为 $z$ 轴，下底面所在平面为 $z=0$。两个大铁球的半径等于圆柱底面半径，
故其球心均在轴线上。圆柱高为 $4$，两个单位球互不重叠且均在容器内，因此它们的球心
只能是
\[
  O_1=(0,0,1),\qquad O_2=(0,0,3).
\]

设一个小钢球的半径为 $r$，球心到圆柱轴线的距离为 $\rho$，球心高度为 $z$。它与两个
大铁球均外切，故
\[
  \rho^2+(z-1)^2=(1+r)^2,
  \qquad
  \rho^2+(z-3)^2=(1+r)^2.
\]
两式相减得 $z=2$，进而
\[
  \rho^2+1=(1+r)^2. \tag{1}
\]
小球位于圆柱内，故 $\rho+r\leqslant1$，特别地 $r\leqslant1$。由于球心高度为 $2$，
小球不可能与上、下底面相切；题设要求它与容器内表面相切，所以它必与侧面相切，即
\[
  \rho=1-r. \tag{2}
\]
联立 (1)、(2)，得到
\[
  r=\frac14,\qquad \rho=\frac34.
\]

所以所有小球的球心均在平面 $z=2$ 内、半径为 $3/4$ 的同一个圆周上。为使小球内部
互不相交，任意两球心的距离至少为 $2r=1/2$。

若有至少 $10$ 个小球，将球心按圆周顺序排列，必有两个相邻球心所对的圆心角不超过
$36^\circ$。这两个球心的距离至多为
\[
  2\cdot\frac34\sin18^\circ
  =\frac{3(\sqrt5-1)}8<\frac12,
\]
矛盾。因此小球至多有 $9$ 个。

另一方面，把九个球心置于该圆的一个内接正九边形的九个顶点上，最近两球心的距离为
\[
  \frac32\sin20^\circ>\frac12.
\]
最后一个不等式可严格验证如下：若 $s=\sin20^\circ\leqslant1/3$，则因
$3s-4s^3$ 在 $[0,1/2]$ 上递增，有
\[
  \frac{\sqrt3}{2}=\sin60^\circ=3s-4s^3
  \leqslant\frac{23}{27}<\frac{\sqrt3}{2},
\]
矛盾。因此九个小球可以放入，最大可能个数为
\[
  \ans{9}.
\]

\section*{11. 最短有趣数列及其个数}

设选出的十项等比子列为
\[
  b_j=a_{i_j}\qquad (j=1,2,\ldots,10),
\]
公比为 $q\ne\pm1$。按等比数列的定义，各项均非零。令
\[
  \Delta_j=b_{j+1}-b_j\qquad(j=1,2,\ldots,9).
\]
原数列相邻两项之差为 $\pm1$，所以任意两项之差均为整数，特别地每个 $\Delta_j$ 都是
非零整数。又由等比关系，
\[
  \Delta_{j+1}=q\Delta_j.
\]
故 $q=\Delta_2/\Delta_1$ 是非零有理数。将其写为最简分数
\[
  q=\frac ps,\qquad p\in\mathbb Z\setminus\{0\},\quad
  s\in\mathbb Z_{>0},\quad \gcd(|p|,s)=1.
\]
因为
\[
  \Delta_9=\Delta_1\frac{p^8}{s^8}\in\mathbb Z,
\]
所以 $s^8\mid\Delta_1$。于是存在非零整数 $D$，使得
\[
  \Delta_j=Dp^{j-1}s^{9-j}\qquad(j=1,2,\ldots,9).
\]

原数列每走一步，数值只改变 $1$，故
\[
  i_{j+1}-i_j\geqslant|b_{j+1}-b_j|=|\Delta_j|.
\]
从而
\[
\begin{aligned}
  m-1&\geqslant i_{10}-i_1\\
     &\geqslant\sum_{j=1}^9|\Delta_j|\\
     &=|D|\bigl(s^8+|p|s^7+\cdots+|p|^8\bigr).
\end{aligned}
\]
由于 $q\ne\pm1$，正整数 $|p|,s$ 不会同时等于 $1$，其中至少一个不小于 $2$。因此
\[
  m-1\geqslant1+2+2^2+\cdots+2^8=511,
\]
即 $m\geqslant512$。

数列
\[
  1,2,3,\ldots,512
\]
含有十项等比子列 $1,2,4,\ldots,512$，所以
\[
  \ans{m_0=512}.
\]

下面计算长度为 $512$ 的有趣数列个数。此时上述各个不等式必须同时取等。因此
\[
  |D|=1,\qquad \{|p|,s\}=\{1,2\},\qquad i_1=1,\qquad i_{10}=512,
\]
并且在每两个相邻的等比子列项之间，原数列必须始终沿同一方向逐次改变 $1$，不能折返。
公比只有
\[
  q\in\left\{2,\frac12,-2,-\frac12\right\}.
\]
由
\[
  \Delta_1=(q-1)b_1=Ds^8
\]
可知，对于每个公比，$D=\pm1$ 分别唯一确定首项：

\begin{center}
\renewcommand{\arraystretch}{1.35}
\begin{tabular}{c c}
\toprule
$q$ & $b_1=a_1$ 的两个可能值\\
\midrule
$2$ & $\pm1$\\
$1/2$ & $\pm512$\\
$-2$ & $\pm\dfrac13$\\
$-1/2$ & $\pm\dfrac{512}{3}$\\
\bottomrule
\end{tabular}
\end{center}

对表中的每一种选择，先取相应的十项等比数列，再在相邻两项间以步长 $+1$ 或 $-1$
单调补齐；总步数恰为 $511$，故得到一个长度为 $512$ 的有趣数列。八种选择的首项互不
相同，所得数列互不相同；上述取等条件也表明不存在其他数列。因此
\[
  \ans{\text{不同的最短有趣数列共有 }8\text{ 个}}.
\]

\end{document}
